\chapter{Evidence And Results}

\section{Evidence Summary}

This chapter should contain the strongest figures, tables, measurements, examples, or qualitative evidence. Each figure or table should support a clear claim in the surrounding text.

\begin{figure}[htbp]
\centering
\begin{tikzpicture}
\begin{axis}[
  width=0.82\textwidth,
  height=6cm,
  ybar,
  ymin=0,
  ylabel={Score},
  symbolic x coords={Baseline,Variant A,Variant B,Proposed},
  xtick=data,
  nodes near coords,
  nodes near coords align={vertical},
  bar width=18pt,
  axis lines*=left,
  tick style={draw=none},
  ymajorgrids=true,
  grid style={Line}
]
\addplot[fill=Accent!75,draw=Accent] coordinates {
  (Baseline,62)
  (Variant A,71)
  (Variant B,76)
  (Proposed,84)
};
\end{axis}
\end{tikzpicture}
\caption{Example pgfplots chart. Replace with real metrics or generated figure files.}
\label{fig:example-chart}
\end{figure}

\section{Interpretation}

Explain what the reader should learn from \Cref{fig:example-chart}. A good evidence section does not merely display data; it tells the reader why the data changes the conclusion.
